% Source: https://artofproblemsolving.com/community/c3h1643631_vietas_formula
% Stable identity: aops-topic-1643631-q01
Let $\Delta=b^2-4ac$ and let $\delta$ be a number with $\delta^2=\Delta$ (if $\Delta<0$ take $\delta=i\sqrt{-\Delta}$). The quadratic formula gives
\[
r_1=\dfrac{-b+\delta}{2a},\qquad r_2=\dfrac{-b-\delta}{2a}.
\]
Then
\[
r_1+r_2=\dfrac{-2b}{2a}=-\dfrac{b}{a},\qquad r_1r_2=\dfrac{(-b)^2-\delta^2}{4a^2}=\dfrac{b^2-(b^2-4ac)}{4a^2}=\dfrac{4ac}{4a^2}=\dfrac{c}{a}.
\]
Consequently, for every $x$,
\[
a(x-r_1)(x-r_2)=a\left(x^2-(r_1+r_2)x+r_1r_2\right)=a\left(x^2+\dfrac{b}{a}x+\dfrac{c}{a}\right)=ax^2+bx+c. \qquad\blacksquare
\]
Nothing here needs the roots to be integers or rational: ``factorable'' only matters if we insist on integer coefficients in the factors. For example, $x^2-x-1$ has roots $\dfrac{1\pm\sqrt5}{2}$, with sum $1$ and product $\dfrac{1-5}{4}=-1$, and $x^2+x+1$ has roots $\dfrac{-1\pm i\sqrt3}{2}$, with sum $-1$ and product $\dfrac{1+3}{4}=1$, exactly as Vieta's formulas predict.
